\documentclass[10pt,twocolumn]{article}
\usepackage[margin=0.8in]{geometry}
\usepackage{booktabs}
\usepackage{amsmath}
\usepackage[hidelinks]{hyperref}
\usepackage{xcolor}
\usepackage{enumitem}
\setlist{nosep,leftmargin=*}

\title{Implied Facts and Cheap Conflict Checks:\\Keeping Long-Term Memory Current with a Small Local Model}
\author{Ali Ashter\\Independent researcher, Calgary, Canada\\\texttt{hello@ashter.dev}}
\date{October 2026}

\begin{document}
\maketitle

\begin{abstract}
Memory layers for LLM assistants are good at recalling what a user said and poor at noticing when a later
conversation quietly overturns it. On STALE, a benchmark of 400 such implicit and explicit conflicts, published
memory systems score 5 to 18\%. We describe \emph{tidemark}, a small memory layer built around three ideas: record
what each message implies as well as what it states, so that implicit conflicts become explicit; check new facts
against old ones with a two-stage test, a one-token probability filter followed by brief reasoning only for flagged
pairs; and always place the user's full current state, with what is no longer true, in the prompt. With an open 8B
model (Qwen3-8B) it scores 55.0 on 300 held-out STALE scenarios, against 8.9 for the same model reading the user's
messages directly. We also find two gaps in how STALE is scored: it rewards forgetting, since it never checks that
still-true facts survive, and it accepts corrections given for the wrong reason. We add a retention probe and a
right-reason check, and report both. Code, prompts and every judged answer are public.
\end{abstract}

\section{Introduction}

A user tells an assistant they have lived in Portland for years. Months later they mention shaking a bark scorpion
out of their boot in the dry heat. Nobody says ``I moved'', yet an assistant that still suggests rainy-day
activities in Portland is wrong. Long-term memory for assistants has mostly been evaluated on recall: finding what the
user said, sometimes across hundreds of sessions \cite{locomo,longmemeval}. Recall is close to solved on those
benchmarks. Keeping memory \emph{current} is not.

STALE \cite{stale} isolates this problem. Each of its 400 scenarios hides an old statement and a later one that
invalidates it, either explicitly (type T1) or only through common sense (type T2), inside about 175K tokens of
conversation. Three probes follow: whether the old state still holds (state resolution, SR), a question that takes
the old state for granted (premise resistance, PR), and a task whose right answer depends on the new state without
naming it (implicit policy adaptation, IPA). Frontier models reading the whole history reach at most 55.2, and
memory systems such as Mem0 \cite{mem0}, Zep \cite{zep}, A-Mem \cite{amem} and LightMem score 5.1 to 17.8. The
authors' own pipeline, CUPMem, reaches 68.0 with GPT-4o-mini and about ten kinds of prompts per session.

We ask how far a small model running locally can get, and what makes the difference. Our contributions:
\begin{itemize}
\item \textbf{Implied facts.} Writing down what a message implies (``the user lives in a hot desert region'') turns
an implicit conflict into one that a similarity search and a cheap check can catch. It is the largest single gain in
our ablation, $+17.4$ points.
\item \textbf{A two-stage conflict check} for small models. One-token yes/no checks are cheap but badly calibrated;
the wording alone moved recall from 1 to 11 of 12 conflicts on a hand-written set. A permissive one-token filter
followed by short reasoning on flagged pairs was both accurate and affordable.
\item \textbf{Full-state readout.} A user's current state is small, so it always goes into the prompt in full,
together with what is no longer true. No retrieval step can miss the fact a question depends on.
\item \textbf{Two scoring fixes for STALE}: a retention probe, without which aggressive forgetting is rewarded, and a
right-reason check, without which corrections justified by unrelated details count as correct.
\end{itemize}

\section{Related work}

\textbf{Memory layers.} MemGPT \cite{memgpt} lets the agent page its own memory. Mem0 \cite{mem0} extracts facts and
asks an LLM whether each new fact should add to, update or delete similar existing ones. Zep \cite{zep} keeps a
temporal knowledge graph whose edges carry validity intervals and are closed when a contradicting edge arrives.
A-Mem \cite{amem} maintains evolving linked notes. All of them find candidates for an update by similarity, which
works for ``I moved to Austin'' against ``I live in Seattle'' but not for a scorpion against Portland.

\textbf{Keeping memory valid.} CUPMem \cite{stale} maintains a structured current-state profile, proposes affected
state ``buckets'' when a practical basis changes, adjudicates invalidations with an LLM and verifies premises at query
time. PACMI \cite{pacmi} tracks typed provenance dependencies and propagates invalidation along them, in the lineage
of truth maintenance systems \cite{doyle}. MemConflict \cite{memconflict} benchmarks memory systems under dynamic,
static and conditional conflicts. Our system shares the write-side invalidation idea with CUPMem but differs in how
candidates are found (implied facts rather than bucket proposals), how pairs are judged (a probability filter before
any generation) and how memory is read (the whole current state rather than retrieved items).

\section{Method}

tidemark keeps, per user, a ledger of short facts. Each fact has a date, a topic, a flag for whether it was stated or
inferred, and, once overturned, a pointer to the fact that overturned it.

\textbf{Only the user's words.} In STALE about 90\% of each conversation is assistant text, which says nothing about
the user. Reading only user turns cuts a scenario from about 175K to about 22K tokens.

\textbf{Extraction with implied facts.} After each conversation the model lists lasting facts about the user's life
as standalone sentences, each with one of 16 topics (home, climate, housing, work, school, family, pets, health, diet,
religion, transport, routine, hobby, activity, travel, request). It then adds at most three facts the messages clearly
imply but never state, marked as inferred. The prompt gives one example: scorpions in one's boots in relentless dry
heat imply living in a hot desert region. Output is constrained to a JSON schema.

\textbf{Conflict check.} Requests and trips never overturn anything. A fact about the current activity (``on a train
for three more hours'') expires when a later activity arrives. Every other new fact is compared with the five most
similar current facts of the same topic group (home, climate and housing form one group) and the three most similar
of any checked topic, by cosine similarity of \texttt{bge-base-en-v1.5} embeddings \cite{bge}. For each pair, stage 1
asks ``taking the later information into account, is the earlier note still true about the user today?'' and reads
$p(\textit{No})$ from the first answer token, renormalised over Yes and No. Pairs with $p(\textit{No})>0.5$ go to stage
2, where the model reasons in thinking mode and returns a JSON verdict; the prompt states that visits, trips, new
interests, unrelated topics and vague guesses overturn nothing. A confirmed verdict marks the old fact as no longer
true.

\textbf{Readout.} The answer prompt contains every current fact, newest first, and every contradicted fact with the
date and the fact that replaced it. Expired activities are left out. The model is told that current facts remain true
and to treat something as changed only if it is on the second list.

\textbf{Premise guard.} Before answering, the five contradicted facts closest to the question are checked with the
same one-token readout: ``does this question assume or rely on the following being true?'' If so, a note telling the
model what changed is appended to the prompt.

\section{Setup}

\textbf{Models.} Answers and memory come from Qwen3-8B \cite{qwen3} in FP8 on vLLM \cite{vllm}, using about 16 GB of
GPU memory. Grading uses Qwen3-14B-AWQ with thinking enabled and STALE's judge prompt unchanged. The paper used
Gemini-3.1-flash-lite, which we could not access, so our numbers and theirs are comparable only roughly. As a check,
our full-history baseline (8.9) lands near the 14.6 reported for Qwen3.5-9B.

\textbf{Split.} All design changes after the first full run were made while inspecting the first 100 scenarios
(interleaved T1 and T2). We report the remaining 300 as held out.

\textbf{Extra metrics.} \emph{Keep}: we ask the model whether the scenario's newer statement is still the case;
nothing later contradicts it, so the answer should be yes. \emph{Right reason}: a second judgment, applied to SR and PR
answers that passed, of whether the correction rests on the actual change rather than an unrelated detail.
\emph{Linkage}: without any judge, whether the memory marked a fact from the old-state session as overturned by a
fact from the new-state session.

\textbf{Baselines.} The same model reading all user messages with dates, once with STALE's neutral prompt and once
with our answer prompt.

\section{Results}

\begin{table*}[t]
\centering
\begin{tabular}{lccccccc}
\toprule
Held-out 300 & SR & PR & IPA & \textbf{Overall} & Right reason & Keep \\
\midrule
User messages, benchmark prompt & 6.7 & 0.7 & 19.3 & 8.9 & 2.5 & 93.0 \\
User messages, our prompt & 15.7 & 1.0 & 21.0 & 12.6 & 6.5 & 97.0 \\
tidemark v1 & 66.0 & 41.7 & 34.0 & 47.2 & 38.4 & 53.3 \\
tidemark, no premise guard & 56.0 & 21.3 & 39.3 & 38.9 & 30.3 & 84.3 \\
\textbf{tidemark} & \textbf{63.7} & \textbf{57.3} & \textbf{44.0} & \textbf{55.0} & \textbf{48.7} & 65.0 \\
\midrule
\multicolumn{7}{l}{\textit{Reported in \cite{stale}, Gemini judge, all 400}} \\
Qwen3.5-9B, full history & & & & 14.6 & & \\
Gemini-3.1-pro, full history & & & & 55.2 & & \\
Mem0 / Zep / A-Mem / LightMem & & & & 8.3 / 6.0 / 5.1 / 17.8 & & \\
CUPMem (GPT-4o-mini) & & & & 68.0 & & \\
\bottomrule
\end{tabular}
\caption{STALE accuracy (\%) on the 300 held-out scenarios. SR, PR and IPA average the T1 and T2 cells; Overall is the
mean of all six cells, as in the benchmark. Right reason (mean of SR and PR) and Keep are our additions.}
\label{tab:main}
\end{table*}

Table~\ref{tab:main} gives the main results. tidemark raises the 8B model from 8.9 to 55.0 overall. Over all 400
scenarios it scores 58.2; on the 100 tuning scenarios 68.0, which overstates the held-out score because those
scenarios are easier (v1, fixed before tuning, also drops from 57.3 to 47.2 between them).

\begin{table}[t]
\centering
\begin{tabular}{lcc}
\toprule
First 100 scenarios, v1 & Overall & Keep \\
\midrule
User messages, our prompt & 11.3 & 97 \\
+ extracted facts & 18.0 & 70 \\
+ two-stage conflict check & 31.3 & 54 \\
+ implied facts & 48.7 & 84 \\
+ premise guard & 57.3 & 48 \\
\bottomrule
\end{tabular}
\caption{What each part adds (v1 pipeline).}
\label{tab:ablation}
\end{table}

\textbf{Ablation.} Table~\ref{tab:ablation} adds one component at a time. Extracted facts alone help little;
conflict checks help more; implied facts give the largest jump ($+17.4$) and also raise retention, because the model
no longer has to guess at changes the memory has made explicit.

\textbf{What the memory caught.} By the linkage measure, v1 tied the right two sessions in 22.5\% of scenarios. The
misses were: the check judged the pair still true or the pair was not among the candidates (27.2\%), old and new facts
were filed under different topic groups (15.8\%), the old fact had already been overturned by something unrelated
(14.5\%), and the old or new fact had a topic that was never checked (18.0\%). The final version, which adds routines
and hobbies to the checked topics, lets a few candidates cross topic groups and lets activities expire, links the right
sessions in 51.0\% of all 400 scenarios, overturns the old state only through unrelated facts in 28.0\% and misses it
in 21.0\%.

\textbf{The premise guard is a trade-off.} It adds 36 points on held-out PR but lowers Keep from 84.3 to 65.0: once
warned, the model is more willing to call things outdated. Without the guard, tidemark scores 38.9 while keeping 84\%
of current facts. Which setting is right depends on whether stale premises or false alarms cost more.

\textbf{Cost.} Per scenario of 50 conversations: 50 extraction calls, about 450 one-token checks and about 170 short
reasoning calls, all of which can run after a conversation ends. At answer time the memory adds about 13K tokens,
against 22K for the raw user messages.

\section{Analysis}

\textbf{One-token checks are badly calibrated.} On 24 hand-written pairs (12 conflicts), the question ``can both be
true today?'' caught 1 conflict, ``is the earlier note still true?'' caught 11 with 1 false alarm, and ``does the later
note make the earlier one outdated?'' caught 1. On 48 harder pairs, including request-style facts and implicit
conflicts, the best wording caught 14 of 16 conflicts but raised 13 false alarms in 32. NLI cross-encoders did no
better (DeBERTa-v3-base: 15/16 with 19 false alarms; large: 13/16 with 11). Requiring the LLM and NLI to agree gave
11/16 with 5. Short reasoning on the pairs the one-token filter flagged caught 13 of 16 with 1 false alarm, at about
380 generated tokens per flagged pair. Since nearly every pair in a real ledger is compatible, the false-alarm rate is
what decides whether a memory keeps working; an early pilot with the one-token check alone kept only 20 to 30\% of
facts current.

\textbf{STALE rewards forgetting.} A memory that marks everything outdated passes most SR and many PR probes. Our
early pilots showed exactly this before the retention probe was added. We suggest that results on STALE report a
retention figure next to the benchmark score.

\textbf{STALE accepts the wrong reason.} The judge passes ``the user may no longer be an independent voter'' even when
the stated reason is that the user later asked about shoe cleaner. In our hand check such answers were common in early
versions. Counting only corrections that cite the actual change, tidemark's held-out SR falls from 63.7 to 50.0 and its PR
from 57.3 to 47.3.

\section{Limitations}

We test one answer model and one judge, and the judge is not the benchmark's. The right-reason check is a second
model judgment, not a human one; the main judge was spot-checked by hand on 12 scenarios. The retention probe reuses
the scenario's new-state message, which sometimes describes a moment, so a ``no longer the case'' answer is
occasionally fair. STALE's filler sessions come from many different people, so some of the ``unrelated'' overturns are
defensible readings of contradictory filler. The memory goes into the prompt in full, which suits dozens of
conversations but would need trimming for years of history. Finally, tidemark keeps the user's situation, not what was
said; it is a complement to transcript retrieval, not a replacement.

\section{Conclusion}

A small local model can keep long-term memory current far better than today's memory layers if it writes down what
messages imply, checks conflicts cheaply but not naively, and always sees the whole current state. Two changes to how
the benchmark is scored, a retention probe and a right-reason check, make the remaining trade-offs visible. Code,
prompts, hand-written test pairs and all judged answers are at \url{https://github.com/alleyxd/tidemark}.

\begin{thebibliography}{99}\small
\bibitem{stale} H.~Chao, Y.~Bai, R.~Sheng, T.~Li, Y.~Sun. STALE: Can LLM agents know when their memories are no longer
valid? arXiv:2605.06527, 2026.
\bibitem{mem0} P.~Chhikara et al. Mem0: Building production-ready AI agents with scalable long-term memory.
arXiv:2504.19413, 2025.
\bibitem{zep} P.~Rasmussen et al. Zep: A temporal knowledge graph architecture for agent memory. arXiv:2501.13956, 2025.
\bibitem{amem} W.~Xu et al. A-MEM: Agentic memory for LLM agents. arXiv:2502.12110, 2025.
\bibitem{memgpt} C.~Packer et al. MemGPT: Towards LLMs as operating systems. arXiv:2310.08560, 2023.
\bibitem{locomo} A.~Maharana et al. Evaluating very long-term conversational memory of LLM agents. ACL, 2024.
\bibitem{longmemeval} D.~Wu et al. LongMemEval: Benchmarking chat assistants on long-term interactive memory. ICLR, 2025.
\bibitem{pacmi} PACMI: Provenance-aware cascading memory invalidation for long-term LLM agents. arXiv:2610.05732, 2026.
\bibitem{memconflict} MemConflict: Evaluating long-term memory systems under memory conflicts. arXiv:2605.20926, 2026.
\bibitem{doyle} J.~Doyle. A truth maintenance system. Artificial Intelligence 12(3), 1979.
\bibitem{qwen3} Qwen Team. Qwen3 technical report. arXiv:2505.09388, 2025.
\bibitem{vllm} W.~Kwon et al. Efficient memory management for large language model serving with PagedAttention. SOSP, 2023.
\bibitem{bge} S.~Xiao et al. C-Pack: Packaged resources to advance general Chinese embedding. arXiv:2309.07597, 2023.
\end{thebibliography}

\end{document}
